% Sammenligning mellom land

\begin{frame}{Landene imellom: hvor Norge havner}
\footnotesize
\setlength{\tabcolsep}{5pt}
\begin{tabular}{l p{5.2cm} p{5.6cm}}
\toprule
 & Beskrivende (casestudier) & Formell estimering (unge, Q5 mot Q1) \\
\midrule
USA     & programvareutviklere 22--25: $-20\,\%$ & $-15$ log-punkter (foretaksfaste effekter) \\
Sverige & $\approx -44\,\%$ rå, øverste kvartil  & $-5{,}5\,\%$ innad i foretak innen 2025H1 \\
Danmark & bred nedgang for unge                  & null; utelukker lønnseffekter over 2\,\% \\
Finland & null                                   & null \\
\midrule
Norge   & programvareutviklere 21--30: $-18\,\%$ & $\approx 0$ i 21--30; 22--25 skiller lag først fra våren 2025 \\
\bottomrule
\end{tabular}

\vspace{0.8em}
\begin{itemize}\setlength{\itemsep}{4pt}
  \item I casestudiene ligner Norge på USA og Sverige.
  \item I den systematiske analysen av hele fordelingen ligner Norge på Danmark og Finland.
  \item Forskjellene mellom land kan gjenspeile ulike forstyrrelser (renteøkninger, normalisering etter pandemien) snarere enn ulike KI-effekter.
\end{itemize}
\end{frame}
